\documentclass[11pt,a4paper]{article}

% ---------------------------------------------------------------
% Langue
% ---------------------------------------------------------------
\usepackage[utf8]{inputenc}
\usepackage[T1]{fontenc}
\usepackage[french]{babel}
\usepackage{lmodern}

% ---------------------------------------------------------------
% Core mathematics packages
% ---------------------------------------------------------------
\usepackage{amsmath, amssymb, amsthm, mathtools}
\usepackage{bm}

% ---------------------------------------------------------------
% Algorithms
% ---------------------------------------------------------------
\usepackage[ruled,vlined,linesnumbered,french]{algorithm2e}

% ---------------------------------------------------------------
% Graphics, figures, and plots
% ---------------------------------------------------------------
\usepackage{graphicx}
\usepackage{tikz}
\usepackage{pgfplots}
\pgfplotsset{compat=1.18}
\usepackage{subcaption}
\usepackage{adjustbox}

% ---------------------------------------------------------------
% Tables
% ---------------------------------------------------------------
\usepackage{booktabs}
\usepackage{tabularx}
\usepackage{siunitx}
\sisetup{output-decimal-marker={,}, retain-explicit-plus}

% ---------------------------------------------------------------
% Cross-referencing and citations
% ---------------------------------------------------------------
\usepackage[numbers,sort&compress]{natbib}
\usepackage[colorlinks=true,linkcolor=blue,citecolor=blue,urlcolor=blue]{hyperref}
\usepackage[capitalise,noabbrev,french]{cleveref}

% ---------------------------------------------------------------
% Typography
% ---------------------------------------------------------------
\usepackage{microtype}
\usepackage{parskip}
\setlength{\parindent}{0pt}
\setlength{\parskip}{6pt}

% ---------------------------------------------------------------
% Page layout
% ---------------------------------------------------------------
\usepackage[margin=2.5cm]{geometry}

% ---------------------------------------------------------------
% Miscellaneous
% ---------------------------------------------------------------
\usepackage{xcolor}
\usepackage{todonotes}

% ---------------------------------------------------------------
% Theorem environments (amsthm)
% ---------------------------------------------------------------
\theoremstyle{definition}
\newtheorem{definition}{Définition}[section]
\newtheorem{hypothese}{Hypothèse}
\renewcommand{\thehypothese}{H\arabic{hypothese}}
\theoremstyle{remark}
\newtheorem{remark}{Remarque}[section]

% ---------------------------------------------------------------
% Notation macros
% ---------------------------------------------------------------
\newcommand{\norm}[1]{\left\lVert #1 \right\rVert}
\newcommand{\ip}[2]{\left\langle #1,\, #2 \right\rangle}
\newcommand{\abs}[1]{\left\lvert #1 \right\rvert}
\newcommand{\R}{\mathbb{R}}
\newcommand{\E}{\mathbb{E}}
\newcommand{\indicator}{\mathbf{1}}
\newcommand{\bigO}[1]{\mathcal{O}\!\left(#1\right)}
\DeclareMathOperator*{\argmin}{arg\,min}
\DeclareMathOperator*{\argmax}{arg\,max}
\DeclareMathOperator{\softmax}{softmax}
\DeclareMathOperator{\CE}{CE}
\DeclareMathOperator{\KL}{KL}
\DeclareMathOperator{\SupCon}{SupCon}
\DeclareMathOperator{\ReLU}{ReLU}
\newcommand{\Src}{\mathcal{S}}
\newcommand{\Tgt}{\mathcal{T}}
\newcommand{\Loss}{\mathcal{L}}
\newcommand{\Wass}{W_{\varepsilon}}

\title{Étiquetage géométrique de la cible par barycentres de Wasserstein\\
  pour l'adaptation de domaine multi-sources et multimodale}
\author{Kevin Djoumessi\\
  Faculté des Sciences, Université de Yaoundé I, Yaoundé, Cameroun\\
  \texttt{kevin.djoumessi@facsciences-uy1.cm}\\[4pt]
  \todo[inline]{Compléter la liste des auteurs et des affiliations.}}
\date{\today}

\begin{document}

\maketitle

\begin{abstract}
Nous étudions l'adaptation de domaine multi-sources et multimodale : plusieurs domaines sources étiquetés, un
domaine cible sans étiquette, deux modalités par exemple. Les méthodes de référence, telles que M2CAN, alignent les
représentations puis entraînent un classifieur sur l'espace fusionné. Nous proposons de retirer ce classifieur et de
lire les classes de la cible dans la géométrie de l'espace fusionné : chaque classe source est résumée par le
barycentre de Wasserstein de ses nuages dans les différentes sources, la cible est partitionnée par un algorithme
spectral, et un transport optimal entre ces structures, suivi d'un appariement hongrois, affecte une classe à chaque
groupe. Cet étiquetage, dérivé de HOT-DA et de WBHR, est exécuté périodiquement pendant l'entraînement et fournit des
pseudo-étiquettes à une perte contrastive supervisée sur l'espace fusionné. Sur des caractéristiques gelées (CLIP,
AST), avec un protocole apparié par cible et par graine et des protocoles écrits avant les mesures, la méthode
atteint le niveau de M2CAN réglé sur deux familles de jeux sur trois (reconnaissance d'actions HAC : 88,64 contre
88,36 ; évaluation esthétique : 65,07 contre 63,99 ; écarts non significatifs), et les pseudo-étiquettes apportent
+7,19 points sur HAC (IC95 $[+2{,}22 ; +12{,}17]$). Sur l'analyse de sentiment, la méthode est inférieure à M2CAN
($-8{,}45$) et passe sous un plancher sans adaptation. Nos diagnostics montrent que l'appariement des groupes aux
classes est optimal dans la quasi-totalité des runs et que le facteur limitant est la partition de la cible, c'est-à-dire
l'existence d'une structure en groupes dans l'espace fusionné. Nous rapportons également deux hypothèses sur le rôle du
barycentre et du transport hiérarchique, et une variante fidèle à WBHR.
\end{abstract}

\section{Introduction}\label{sec-intro}

L'adaptation de domaine non supervisée cherche à prédire les étiquettes d'un domaine cible, pour lequel seules des
entrées sont disponibles, à partir d'un ou plusieurs domaines sources étiquetés dont la distribution diffère
\citep{courty2017otda}. Le cadre multi-sources et multimodal (MSMMDA) ajoute deux difficultés : les sources sont
plusieurs et ne se recouvrent pas, et chaque exemple est décrit par plusieurs modalités, par exemple un texte et une
image, ou une piste audio et une vidéo \citep{zhao2025msmmda, dong2026mmadaptsurvey}. La méthode de référence pour ce
cadre, M2CAN \citep{zhao2025msmmda}, encode chaque modalité, aligne les représentations par des pertes contrastives
entre modalités et entre domaines, les fusionne par un produit bilinéaire de rang faible \citep{kim2017mlb}, puis
entraîne un classifieur sur l'espace fusionné.

Le transport optimal offre une autre lecture de l'adaptation. HOT-DA \citep{elhamri2022hotda,
elhamri2025theoretical} représente chaque domaine comme une mesure sur des structures, les classes pour la source et
des groupes pour la cible, et résout un problème de transport entre structures avant d'apparier chaque groupe à une
classe. WBHR \citep{elhamri2025wbhr} étend cette représentation hiérarchique au cas multi-sources en résumant chaque
classe, à travers les sources, par un barycentre de Wasserstein \citep{cuturi2014fast, montesuma2021wbtransport}.
Ces méthodes opèrent sur des caractéristiques fixes, après coup.

Nous proposons de placer cet étiquetage géométrique à l'intérieur de l'entraînement d'une fusion multimodale, et de
supprimer le classifieur sur l'espace fusionné. Les classes de la cible sont lues par la géométrie : barycentres de
Wasserstein des classes sources, partition spectrale de la cible, transport entre structures et appariement. Cet
étiquetage est exécuté périodiquement pendant l'apprentissage ; les étiquettes qu'il produit, filtrées par l'accord des
têtes unimodales, entrent dans une perte contrastive supervisée sur l'espace fusionné. L'espace que la géométrie lit
est ainsi façonné par ce qu'elle y lit.

Nos contributions sont les suivantes.
\begin{enumerate}
  \item Une méthode d'adaptation multi-sources et multimodale sans classifieur sur l'espace fusionné, où un étiquetage
    géométrique dérivé de WBHR et de HOT-DA sert d'oracle pendant l'entraînement (\cref{sec-methode}). La méthode
    accepte un nombre quelconque de sources et de modalités.
  \item Une comparaison appariée à M2CAN sur trois familles de jeux (actions, esthétique, sentiment), à caractéristiques
    gelées identiques, mêmes graines et mêmes découpes, avec des protocoles écrits avant les mesures
    (\cref{sec-protocole,sec-resultats}). La méthode atteint le niveau de M2CAN sur deux familles et lui est inférieure
    sur la troisième.
  \item Une analyse du mécanisme : l'apport des pseudo-étiquettes, deux hypothèses sur le barycentre et sur le
    transport hiérarchique, la variante fidèle à WBHR, et des diagnostics qui localisent le facteur limitant dans la
    partition de la cible plutôt que dans l'appariement (\cref{sec-analyse}).
\end{enumerate}

Le code, les configurations et les protocoles sont publics.\footnote{\url{https://github.com/Kevindjoum/MM-WBHR}}

\section{Travaux liés}\label{sec-lies}

\paragraph{Adaptation multi-sources et multimodale.}
M2CAN \citep{zhao2025msmmda} définit le cadre MSMMDA et propose six jeux de données textes et images (sentiment :
TumEmo, Twitter, Yelp ; esthétique : AVA, PCCD, RPCD). Sa fonction de coût combine une entropie croisée sur des têtes
unimodales, un alignement contrastif entre modalités pondéré par la cohérence des têtes (CMCFA), un alignement
contrastif supervisé entre domaines (CDCFA) et un classifieur sur la fusion MLB. Nous reprenons les trois premiers termes
et remplaçons le quatrième. SimMMDG \citep{dong2023simmmdg} introduit le jeu HAC (humains, animaux, dessins animés)
pour la généralisation multimodale ; nous l'utilisons en adaptation. Les revues de \citet{dong2026mmadaptsurvey} et
\citet{li2025mmalignsurvey} situent ces travaux.

\paragraph{Transport optimal pour l'adaptation.}
\citet{courty2017otda} transportent les sources vers la cible point à point. HOT-DA \citep{elhamri2022hotda} transporte
des structures (classes sources, groupes cibles) et apparie chaque groupe à une classe ; \citet{elhamri2025theoretical}
en donnent des garanties. WBHR \citep{elhamri2025wbhr} résume chaque classe multi-sources par un barycentre de
Wasserstein à support libre \citep{cuturi2014fast} avant d'appliquer HOT-DA, puis entraîne un classifieur sur les
représentants transportés. \citet{montesuma2021wbtransport} utilisent aussi un barycentre de Wasserstein pour agréger
des sources. Les calculs reposent sur la régularisation entropique \citep{cuturi2013sinkhorn, peyre2019computationalot}.

\paragraph{Pseudo-étiquettes et maximisation de l'information.}
SHOT \citep{liang2020shot} combine une perte de maximisation de l'information sur la cible et des pseudo-étiquettes
par centres de classes. Nos pseudo-étiquettes viennent d'une partition de la cible appariée aux classes sources, et
nous évaluons la maximisation de l'information comme option (\cref{sec-analyse}).

\section{Méthode}\label{sec-methode}

\subsection{Problème et notations}

On dispose de $N$ domaines sources étiquetés $\Src_1, \dots, \Src_N$ et d'un domaine cible $\Tgt$ sans étiquette.
Chaque exemple possède $M$ modalités ordonnées $x = (x_0, \dots, x_{M-1})$ et appartient à l'une de $K$ classes. Le
train cible sert à l'entraînement sans ses étiquettes ; le test cible sert à l'évaluation. La méthode est transductive :
elle partitionne un ensemble de points cibles et ne prédit pas un point isolé.

\subsection{Modèle}

\paragraph{Encodeurs et têtes unimodales.}
Pour la modalité $m$, une dorsale gelée $b_m$ fournit un vecteur de caractéristiques ; un extracteur linéaire
$W_m \in \R^{256 \times d_m}$ suivi d'une normalisation donne
\begin{equation}\label{eq-encodeur}
  e_m = \frac{\ReLU(W_m\, b_m(x_m))}{\norm{\ReLU(W_m\, b_m(x_m))}} \in \R^{256}.
\end{equation}
Une tête unimodale $g_m$, perceptron $256 \to 128 \to 64 \to K$, donne la prédiction
$p_m = \softmax(g_m(e_m))$.

\paragraph{Fusion.}
Les représentations sont fusionnées par un produit bilinéaire de rang faible \citep{kim2017mlb}, le même que M2CAN,
vers un espace de dimension $256$ :
\begin{equation}\label{eq-mlb}
  z = P\Bigl(\sigma\bigl(\textstyle\bigodot_{m} U_m e_m\bigr)\Bigr) + \sum_{m} h_m(e_m),
\end{equation}
avec $U_m \colon \R^{256} \to \R^{128}$, $P \colon \R^{128} \to \R^{256}$, $h_m \colon \R^{256} \to \R^{256}$ linéaires,
$\sigma$ la sigmoïde et $\odot$ le produit terme à terme. Aucun classifieur n'est entraîné sur $z$ : les classes de la
cible y sont lues par la procédure de la \cref{sec-etiquetage}.

\subsection{Outils de transport optimal}\label{sec-ot}

Le coût entre deux vecteurs est le coût cosinus
$c(u, v) = \norm{u / \norm{u} - v / \norm{v}}^2 = 2 - 2\cos(u, v) \in [0, 4]$. Pour une matrice de coût $C$ et deux
mesures discrètes $a$ et $b$, le plan de transport régularisé par l'entropie
\begin{equation}\label{eq-sinkhorn}
  \Gamma^\star = \argmin_{\Gamma \indicator = a,\ \Gamma^\top \indicator = b}\ \ip{\Gamma}{C} - \varepsilon H(\Gamma)
\end{equation}
est calculé par l'algorithme de Sinkhorn en domaine logarithmique \citep{cuturi2013sinkhorn}. Le coût de transport
entre deux nuages $X$ et $Y$ de poids uniformes est $\Wass(X, Y) = \ip{\Gamma^\star}{C}$ avec $C_{ij} = c(x_i, y_j)$.

Le barycentre de Wasserstein à support libre de $N$ nuages $X_1, \dots, X_N$ \citep{cuturi2014fast} est un nuage
$B = (b_1, \dots, b_S)$ de poids $1/S$ minimisant $\sum_k \Wass(B, X_k)$. Il est calculé par point fixe : $B$ est
initialisé par $S$ points tirés dans l'union des nuages ; à chaque itération, le plan $\Gamma_k$ entre $B$ et $X_k$ est
calculé pour chaque $k$, puis chaque point est remplacé par la moyenne de ses projections barycentriques,
\begin{equation}\label{eq-barycentre}
  b_s \leftarrow \frac{1}{N} \sum_{k=1}^{N} S \sum_{j} \Gamma_k[s, j]\, x_{k, j}.
\end{equation}
Nous utilisons $S = 16$, $\varepsilon = 0{,}02$, $60$ itérations de Sinkhorn et $8$ itérations de point fixe, comme
l'algorithme 2 de WBHR \citep{elhamri2025wbhr}. La \cref{sec-analyse} montre que la valeur de $S$ n'a pas d'effet.

\subsection{Étiquetage géométrique de la cible}\label{sec-etiquetage}

La procédure \textsc{Étiqueter} (\cref{alg-etiqueter}) reçoit les représentations $z$ et les étiquettes des sources,
et les $z$ de $n$ points cibles. Elle renvoie une classe pour chaque point cible.

\begin{algorithm}[ht]
\caption{\textsc{Étiqueter} : étiquetage géométrique de la cible}\label{alg-etiqueter}
\KwIn{représentations et étiquettes des sources $\{(z, y)\}_{k \le N}$ ; représentations cibles $Z_\Tgt \in \R^{n \times 256}$ ; $K$}
\KwOut{étiquettes $\hat y \in \{1, \dots, K\}^n$, groupes, marges}
\ForEach{classe $h$}{
  tirer au plus $400$ points de classe $h$ dans chaque source\;
  $B_h \leftarrow$ barycentre de Wasserstein de ces $N$ nuages, \cref{eq-barycentre} \tcp*{représentant de la classe}
}
$C_{ij} \leftarrow c(z_i, z_j)$ pour $z_i, z_j \in Z_\Tgt$ ; $\sigma \leftarrow$ médiane des $C_{ij} > 0$\;
$A \leftarrow \exp(-C / \sigma)$, $A_{ii} \leftarrow 1$\;
$G_1, \dots, G_K \leftarrow$ partition spectrale de $A$ en $K$ groupes \citep{vonluxburg2007spectral}\;
\ForEach{classe $h$, groupe $l$}{
  $W_{hl} \leftarrow \Wass(B_h, G_l)$ \tcp*{coût entre structures}
}
$\Gamma \leftarrow$ plan de Sinkhorn sur $W$, marges uniformes $1/K$, \cref{eq-sinkhorn}\;
$\pi \leftarrow$ appariement hongrois \citep{kuhn1955hungarian} maximisant $\sum_l \Gamma_{\pi(l), l}$\;
$\hat y_j \leftarrow \pi(\text{groupe}(j))$ pour tout $j$\;
$\text{marge}_j \leftarrow \min_{l \ne \text{groupe}(j)} c(z_j, \bar c_l) \,/\, c(z_j, \bar c_{\text{groupe}(j)})$, avec $\bar c_l$ le centre du groupe $l$\;
\Return{$\hat y$, groupes, marges}
\end{algorithm}

Les étapes 4 à 10 sont celles de HOT-DA \citep{elhamri2022hotda} ; le représentant de classe par barycentre
multi-sources est celui de WBHR \citep{elhamri2025wbhr}. Deux différences sont à noter. L'appariement est hongrois,
donc bijectif, là où le texte de WBHR prend un $\argmax$ par classe qui peut envoyer deux classes sur le même groupe.
Et l'étiquette finale est celle du groupe apparié : WBHR transporte au contraire les représentants vers leur groupe et
entraîne un classifieur sur les points transportés. Cette variante est évaluée à la \cref{sec-analyse}.

\subsection{Phase d'entraînement}\label{sec-phase1}

À chaque pas, on tire quatre exemples dans chaque domaine, sources puis cible, avec une vue faible $x$ et une vue
augmentée $\tilde x$. La perte est
\begin{equation}\label{eq-perte}
  \Loss = \Loss_{\text{têtes}} + \Loss_{\text{CMCFA}} + \Loss_{\text{CDCFA}} + \Loss^{\text{fus}}_{\text{SC}},
\end{equation}
avec des poids égaux à $1$. Les trois premiers termes sont ceux de M2CAN \citep{zhao2025msmmda}.

$\Loss_{\text{têtes}}$ est l'entropie croisée des têtes unimodales sur les sources,
$\frac{1}{MN} \sum_{i \le N} \sum_m \CE(g_m(e_m^{(i)}), y^{(i)})$.

$\Loss_{\text{CMCFA}}$ aligne les modalités de chaque domaine, cible comprise, sans étiquette. Pour un exemple $x$ d'un
lot, avec $a = e_0$ et $b = e_1$, l'InfoNCE sur les quatre combinaisons de vues vaut
\begin{equation}\label{eq-infonce}
  \ell(x) = -\log \frac{\sum_{(u, v)} \exp\ip{u_x}{v_x}}{\sum_{(u, v)} \sum_{x'} \exp\ip{u_x}{v_{x'}}},
  \qquad (u, v) \in \{a, \tilde a\} \times \{b, \tilde b\},
\end{equation}
et la cohérence entre les têtes est $V(x) = \KL(p_1 \,\Vert\, p_0)$ pour $M = 2$ (divergence de Jensen-Shannon des $M$
têtes au-delà). Les exemples où les modalités se contredisent pèsent moins :
$\Loss_{\text{CMCFA}} = \sum_{d} \bigl[\tfrac14 \sum_{x} \ell(x) e^{-V(x)} + \tfrac14 \sum_{x} V(x)\bigr]$. Avec
$M \ge 3$, l'InfoNCE est moyenné sur les paires de modalités.

$\Loss_{\text{CDCFA}}$ est une perte contrastive supervisée \citep{khosla2020supcon} par modalité, sur les exemples
sources réunis, de température $\tau = 0{,}07$ :
\begin{equation}\label{eq-supcon}
  \SupCon(Z, Y, \tau) = -\frac{1}{\abs{A}} \sum_{a \in A} \frac{1}{\abs{P(a)}} \sum_{p \in P(a)}
  \log \frac{\exp(z_a \cdot z_p / \tau)}{\sum_{b \ne a} \exp(z_a \cdot z_b / \tau)},
\end{equation}
où $P(a)$ réunit les exemples de même étiquette que $a$ et $A$ ceux qui en ont au moins un.

$\Loss^{\text{fus}}_{\text{SC}}$ est le terme propre à la méthode : une perte $\SupCon$ sur l'espace fusionné $z$,
portant sur les exemples sources et sur les exemples cibles du lot qui possèdent une pseudo-étiquette,
\begin{equation}\label{eq-scfus}
  \Loss^{\text{fus}}_{\text{SC}} = \SupCon\Bigl(\{z^{(i)}\}_{i \le N} \cup \{z^{\Tgt}_j\}_{j \in \mathcal{P}},\
  \{y^{(i)}\}_{i \le N} \cup \{\hat y_j\}_{j \in \mathcal{P}},\ \tau\Bigr).
\end{equation}
C'est le seul terme où les pseudo-étiquettes agissent : un point cible étiqueté $h$ est attiré vers les exemples sources
de la classe $h$ dans l'espace fusionné.

\subsection{Pseudo-étiquettes}\label{sec-pseudo}

Tous les $750$ pas, soit $15$ tours $r = 0, \dots, 14$ sur $3$ époques, le modèle passe en évaluation. Les $z$ des
sources et ceux d'au plus $3000$ exemples du train cible sont calculés, ainsi que la prédiction moyenne des têtes
$\bar p_j = \frac{1}{M} \sum_m p_{m, j}$. La procédure \textsc{Étiqueter} fournit $\hat y_j$, le groupe et la marge de
chaque point cible. Deux filtres s'appliquent.
\begin{enumerate}
  \item \emph{Accord.} Un point n'est retenu que si la prédiction moyenne des têtes désigne la même classe que la
    géométrie, $\argmax_h \bar p_{j, h} = \hat y_j$. Les têtes, fusion tardive des décisions unimodales, ne décident
    pas ; elles filtrent la décision géométrique, fusion intermédiaire des représentations.
  \item \emph{Quota croissant.} Avec $\rho_r = (0{,}3 ; 0{,}6 ; 0{,}9 ; 1{,}0)[\min(r, 3)]$ et $n$ points étiquetés, on
    garde dans chaque classe au plus $\lfloor \rho_r n / K \rfloor$ points, ceux de plus grande marge.
\end{enumerate}
La table $\{\text{identifiant} \mapsto \hat y_j\}$ des points retenus remplace celle du tour précédent et sert jusqu'au
suivant. Les vraies étiquettes du train cible ne servent qu'à journaliser la justesse des pseudo-étiquettes.

\subsection{Évaluation}

Après chaque tour, \textsc{Étiqueter} est appliquée au test cible et l'exactitude des $\hat y$ est journalisée, ainsi
que deux diagnostics : l'\emph{oracle}, exactitude obtenue avec le meilleur appariement possible des groupes aux
classes, qui mesure la qualité de la seule partition ; et le \emph{plancher logistique}, exactitude d'une régression
logistique entraînée sur les $z$ des sources et appliquée au test cible, qui mesure ce que vaut $z$ sans étiquetage
géométrique. Le résultat rapporté est celui du \emph{dernier} point d'évaluation. Le meilleur point, choisi avec les
étiquettes du test, serait optimiste ; nous ne le rapportons jamais.

\paragraph{Hyperparamètres.}
Adam, $\beta = (0{,}9 ; 0{,}999)$, taux $5 \times 10^{-5}$ pour les extracteurs, la fusion et les têtes, décroissance
polynomiale $(1 - t/T)^{0{,}9}$ sur $T = 3 \times 3750$ pas ; lot de $4$ exemples par domaine ; $e_m$ et $z$ en
dimension $256$ ; $\tau = 0{,}07$ ; barycentre : $S = 16$, $\varepsilon = 0{,}02$, $60 \times 8$ itérations, $400$
points par classe et par source ; transport entre structures : $\varepsilon = 0{,}02$, $300$ itérations ;
pseudo-étiquettes : période $750$, $3000$ points, filtre d'accord, quotas $(0{,}3 ; 0{,}6 ; 0{,}9 ; 1)$.

\section{Protocole expérimental}\label{sec-protocole}

\paragraph{Jeux de données.}
Trois familles, chacune de trois domaines qui servent de cible à tour de rôle, les deux autres étant les sources
(\cref{tab-jeux}). HAC \citep{dong2023simmmdg} comprend des vidéos d'actions humaines, animales et de dessins
animés, avec audio et vidéo. Les familles sentiment et esthétique sont celles de M2CAN \citep{zhao2025msmmda}, avec
texte et image. Les découpes sont celles des auteurs de M2CAN et de SimMMDG.

\begin{table}[ht]
\centering
\caption{Jeux de données. Le test cible est l'ensemble sur lequel l'exactitude est mesurée.}\label{tab-jeux}
\small\begin{tabular}{l l l c c}
\toprule
famille & domaines & modalités (dimension) & $K$ & test cible \\
\midrule
HAC & human, animal, cartoon & audio AST (768), vidéo CLIP (512) & 7 & 276 / 176 / 218 \\
esthétique & AVA, PCCD, RPCD & texte CLIP (512), image CLIP (512) & 2 & 847 \\
sentiment & TumEmo, Twitter, Yelp & texte CLIP (512), image CLIP (512) & 3 & 1500 \\
\bottomrule
\end{tabular}
\end{table}

\paragraph{Caractéristiques gelées.}
Toutes les méthodes reçoivent les mêmes vecteurs, extraits une fois par des encodeurs gelés : CLIP ViT-B/32
\citep{radford2021clip} pour le texte, l'image et la vidéo (moyenne de huit images par segment temporel
\citep{wang2016tsn}), AST \citep{gong2021ast} pour l'audio. Chaque exemple a une vue faible et quatre vues fortes,
avec les augmentations de M2CAN. Pour M2CAN, les vecteurs sont complétés par des zéros jusqu'aux dimensions attendues
par son code.

\paragraph{Méthodes comparées.}
M2CAN est exécuté à partir du code de ses auteurs, adapté aux caractéristiques gelées, avec la configuration publiée de
chaque cible et $10$ époques. Ses taux d'apprentissage publiés sont prévus pour des encodeurs ajustés de bout en bout ;
sur caractéristiques gelées, ils figent presque ses couches entraînables. Nous rapportons donc aussi \emph{M2CAN réglé},
dont le taux d'extracteur est porté à $5 \times 10^{-5}$, le même que celui de notre méthode. C'est la référence
retenue pour toutes les comparaisons. La méthode est entraînée $3$ époques. Nous rapportons également la méthode sans
pseudo-étiquettes (ablation), un plancher \emph{sources seules} (régression logistique sur les deux modalités
concaténées, entraînée sur les sources) et un plancher d'auto-apprentissage naïf (le même modèle avec cinq tours de
pseudo-étiquetage par confiance).

\paragraph{Statistiques.}
Trois graines ($123456$, $234567$, $345678$) par cible, soit $9$ runs par configuration et par famille. Les écarts entre
deux configurations sont appariés par (cible, graine), $9$ paires, avec un intervalle de confiance à $95\,\%$ de Student
à $8$ degrés de liberté. Un écart est déclaré en faveur d'une configuration si l'intervalle exclut zéro et si au moins
$7$ paires sur $9$ sont de son côté ; il est déclaré \emph{non significatif} sinon. Les protocoles (configurations,
comparaisons, critère) ont été écrits avant chaque série de mesures. Les runs sont déterministes : la méthode relancée
avec un code ultérieur reproduit ses chiffres à l'identique.

\section{Résultats}\label{sec-resultats}

\subsection{Comparaison à M2CAN}

\begin{table}[ht]
\centering
\caption{Exactitude au dernier point ($\%$), moyenne sur trois graines, caractéristiques gelées. En gras, la meilleure
  valeur de chaque colonne parmi les méthodes d'adaptation. Les planchers ne sont pas appariés aux autres lignes
  (sous-échantillons différents).}\label{tab-principal}
\small
\begin{tabularx}{\linewidth}{l *{4}{>{\centering\arraybackslash}X}}
\toprule
 & human & animal & cartoon & moyenne \\
\midrule
\multicolumn{5}{l}{\emph{HAC}} \\
méthode & \textbf{93,36} & \textbf{96,40} & 76,15 & \textbf{88,64} \\
méthode sans pseudo-étiquettes & 87,32 & 95,08 & 61,93 & 81,44 \\
M2CAN réglé & 87,68 & 95,45 & \textbf{81,96} & 88,36 \\
M2CAN publié & 86,23 & 86,17 & 77,06 & 83,16 \\
sources seules / auto-apprentissage & 73,79 / 88,04 & 86,55 / 92,80 & 71,10 / 81,35 & 77,15 / 87,40 \\
\midrule
 & AVA & PCCD & RPCD & moyenne \\
\midrule
\multicolumn{5}{l}{\emph{esthétique}} \\
méthode & 72,06 & 62,06 & 61,08 & 65,07 \\
méthode sans pseudo-étiquettes & \textbf{73,04} & \textbf{63,75} & \textbf{61,67} & \textbf{66,16} \\
M2CAN réglé & 69,62 & 62,42 & 59,94 & 63,99 \\
M2CAN publié & 62,38 & 61,59 & 58,24 & 60,74 \\
sources seules / auto-apprentissage & 65,80 / 66,98 & 59,74 / 59,35 & 54,55 / 59,46 & 60,03 / 61,93 \\
\midrule
 & TumEmo & Twitter & Yelp & moyenne \\
\midrule
\multicolumn{5}{l}{\emph{sentiment}} \\
méthode & 46,02 & 56,36 & 39,67 & 47,35 \\
méthode sans pseudo-étiquettes & 44,80 & 56,38 & 44,87 & 48,68 \\
M2CAN réglé & \textbf{58,02} & \textbf{59,07} & \textbf{50,31} & \textbf{55,80} \\
M2CAN publié & 52,18 & 47,89 & 37,73 & 45,93 \\
sources seules / auto-apprentissage & 51,64 / 55,13 & 52,07 / 55,49 & 50,44 / 51,31 & 51,38 / 53,98 \\
\bottomrule
\end{tabularx}
\end{table}

\begin{table}[ht]
\centering
\caption{Écarts appariés par (cible, graine), en points d'exactitude : moyenne, IC95 et nombre de paires en faveur
  de la première configuration. Verdict selon le critère de la \cref{sec-protocole}.}\label{tab-ecarts}
\small
\begin{tabular}{llrcl}
\toprule
famille & comparaison & écart & IC95 & paires / verdict \\
\midrule
HAC & méthode $-$ M2CAN réglé & $+0{,}27$ & $[-3{,}68 ; +4{,}22]$ & 5/9, non significatif \\
HAC & méthode $-$ sans pseudo-étiquettes & $+7{,}19$ & $[+2{,}22 ; +12{,}17]$ & 8/9, \textbf{mieux} \\
HAC & M2CAN réglé $-$ M2CAN publié & $+5{,}21$ & $[+2{,}41 ; +8{,}01]$ & 8/9, mieux \\
esthétique & méthode $-$ M2CAN réglé & $+1{,}08$ & $[-0{,}13 ; +2{,}28]$ & 7/9, non significatif \\
esthétique & méthode $-$ sans pseudo-étiquettes & $-1{,}09$ & $[-1{,}75 ; -0{,}43]$ & 0/9, moins bien \\
esthétique & M2CAN réglé $-$ M2CAN publié & $+3{,}25$ & $[+0{,}86 ; +5{,}64]$ & 9/9, mieux \\
sentiment & méthode $-$ M2CAN réglé & $-8{,}45$ & $[-12{,}32 ; -4{,}58]$ & 0/9, \textbf{moins bien} \\
sentiment & méthode $-$ sans pseudo-étiquettes & $-1{,}33$ & $[-3{,}70 ; +1{,}03]$ & 5/9, non significatif \\
sentiment & M2CAN réglé $-$ M2CAN publié & $+9{,}87$ & $[+6{,}40 ; +13{,}33]$ & 9/9, mieux \\
\bottomrule
\end{tabular}
\end{table}

Les \cref{tab-principal,tab-ecarts} donnent le résultat principal. Sur HAC et sur l'esthétique, la méthode atteint le
niveau de M2CAN réglé : $+0{,}27$ et $+1{,}08$ point, écarts non significatifs. Un étiquetage purement géométrique de
l'espace fusionné y fait donc aussi bien qu'un classifieur entraîné sur cet espace. Sur HAC, la méthode devance M2CAN sur
human ($+5{,}7$) et animal ($+1{,}0$) et reste derrière sur cartoon ($-5{,}8$). Sur le sentiment, la méthode est
inférieure à M2CAN réglé de $8{,}45$ points, sur les neuf paires, et passe sous le plancher sources seules ($47{,}35$
contre $51{,}38$) : l'étiquetage géométrique y perd de l'information que $z$ contient.

Le réglage du taux d'extracteur de M2CAN est déterminant : il apporte $+5$ à $+10$ points selon la famille. Les
comparaisons à M2CAN publié, favorables à la méthode, tiennent surtout à cette asymétrie de réglage et ne sont pas
retenues.

\subsection{Apport des pseudo-étiquettes}

Sur HAC, les pseudo-étiquettes apportent $+7{,}19$ points (IC95 $[+2{,}22 ; +12{,}17]$, $8/9$) ; sur cartoon, la cible
la plus éloignée, elles font passer l'exactitude de $61{,}93$ à $76{,}15$. Les pseudo-étiquettes retenues sont justes à
$94$ à $97\,\%$ en fin d'entraînement sur human et animal, à $79$ à $85\,\%$ sur cartoon. Sur une graine, la divergence
MMD$^2$ entre la cible et les sources dans $z$ est divisée par trois à six (cartoon : $0{,}120 \to 0{,}020$) et
l'information mutuelle normalisée de la partition augmente ($0{,}51 \to 0{,}65$) : les pseudo-étiquettes rapprochent la
cible des classes sources dans l'espace fusionné. Sur l'esthétique et le sentiment, elles n'apportent rien, et coûtent
même $1{,}09$ point sur l'esthétique : la MMD$^2$ n'y bouge pas. Leur bénéfice dépend donc de la justesse de l'étiquetage
géométrique qui les produit.

\section{Analyse du mécanisme}\label{sec-analyse}

\subsection{L'appariement n'est pas le facteur limitant}

Dans $184$ runs sur $198$, l'exactitude au dernier point égale l'oracle de la partition : l'appariement hongrois des
groupes aux classes est optimal. Les exceptions sont concentrées sur le sentiment (écart moyen $0{,}21$ point, maximum
$10{,}5$). Tout ce qui sépare la méthode d'une exactitude parfaite tient donc à la partition de la cible, c'est-à-dire à
l'existence, dans $z$, de $K$ groupes qui coïncident avec les classes. Ce constat éclaire les trois analyses qui
suivent.

\subsection{Hypothèses sur le barycentre et le transport hiérarchique}

\begin{hypothese}\label{hyp-h1}
Résumer une classe par le barycentre de Wasserstein de ses nuages dans les sources représente mieux cette classe que la
concaténation de ces nuages.
\end{hypothese}

\begin{hypothese}\label{hyp-h2}
Le transport hiérarchique (partition, transport entre structures, appariement) étiquette mieux la cible qu'un transport
direct des points cibles vers les points des représentants de classe.
\end{hypothese}

La \cref{tab-h1h2} rapporte les deux hypothèses sur HAC, avec un barycentre dont le support a la taille moyenne de la
classe dans les sources, afin que barycentre et concaténation aient des cardinalités comparables.

\begin{table}[ht]
\centering
\caption{H1 et H2 sur HAC, exactitude au dernier point ($\%$), moyenne sur trois graines.}\label{tab-h1h2}
\begin{tabular}{lcccc}
\toprule
représentant, transport & human & animal & cartoon & moyenne \\
\midrule
barycentre, hiérarchique (méthode) & \textbf{93,36} & \textbf{96,40} & 76,15 & \textbf{88,64} \\
concaténation, hiérarchique & \textbf{93,36} & \textbf{96,40} & 76,15 & \textbf{88,64} \\
barycentre, plat & 91,55 & 90,53 & \textbf{76,91} & 86,33 \\
concaténation, plat & 92,27 & 90,53 & \textbf{76,91} & 86,57 \\
\bottomrule
\end{tabular}
\end{table}

\emph{H1 est partiellement vérifiée.} Le barycentre est plus proche, au sens du coût de transport, de la classe cible
correspondante que la concaténation, dans $9$ cas sur $9$ ($+0{,}022$ en moyenne). Mais il produit exactement les mêmes
étiquettes : le représentant agit sur l'appariement, déjà optimal. Pour la même raison, la taille du support du
barycentre ($16$ points ou taille moyenne de la classe) ne change aucune étiquette sur les trois familles.

\emph{H2 est une tendance favorable, non significative.} Le transport hiérarchique devance le transport plat de
$+2{,}31$ point (IC95 $[-0{,}04 ; +4{,}65]$, $7/9$), surtout sur animal ($+5{,}9$). L'intervalle touche zéro ; le
protocole conclut à un résultat non concluant de justesse.

\subsection{Variante fidèle à WBHR}

WBHR \citep{elhamri2025wbhr} ne s'arrête pas à l'appariement : les points de chaque représentant $B_h$ sont transportés
vers leur groupe apparié par projection barycentrique, $\tilde B_h = \abs{B_h}\, \Gamma^\star_{h, \hat l(h)}\, Z_{\hat l(h)}$,
un classifieur est entraîné sur $\bigcup_h \tilde B_h$ étiqueté $h$ et appliqué à tous les points cibles. Nous avons
mesuré cette fin (appariement par $\argmax$, régression logistique sur les représentants transportés), seule et avec le
support à la taille moyenne de la classe (\cref{tab-wbhr}).

\begin{table}[ht]
\centering
\caption{Variante fidèle à WBHR, exactitude au dernier point ($\%$), moyenne de neuf runs, et écart apparié à la
  méthode.}\label{tab-wbhr}
\small
\begin{tabular}{lccc}
\toprule
 & HAC & esthétique & sentiment \\
\midrule
méthode & 88,64 & 65,07 & 47,35 \\
support à la taille moyenne & 88,64 & 65,07 & 47,35 \\
fin WBHR (argmax, classifieur) & 88,10 & 65,18 & 50,88 \\
les deux & \textbf{89,34} & \textbf{65,16} & \textbf{50,80} \\
\midrule
les deux $-$ méthode & $+0{,}70$ n.s. & $+0{,}09$ n.s. & $+3{,}45$ $[+1{,}28 ; +5{,}63]$, 9/9 \\
les deux $-$ M2CAN réglé & $+0{,}97$ n.s. & $+1{,}17$ n.s. & $-5{,}00$ $[-7{,}56 ; -2{,}44]$ \\
\bottomrule
\end{tabular}
\end{table}

Le classifieur de WBHR n'aide que sur le sentiment ($+3{,}45$, $9/9$), où il dépasse l'oracle de la partition de
$2{,}3$ points : il n'est plus limité par les groupes et corrige une partie des erreurs de la partition. Il ne change
rien là où la partition est bonne, et ne suffit pas à rattraper M2CAN réglé ni le plancher sources seules sur le
sentiment. La méthode garde donc l'étiquetage par groupe, sans classifieur ; la variante reste une option documentée.

\subsection{Leviers sur l'espace fusionné}

Six configurations modifiant un seul élément de la méthode ont été comparées, avec un protocole écrit avant les
mesures (\cref{tab-leviers}) : la maximisation de l'information de SHOT \citep{liang2020shot} sur les têtes du lot
cible (poids $0{,}3$ et $1$), la perte $\SupCon$ de la fusion sur $z$ normalisé, une partition spectrale sur le graphe
des $15$ plus proches voisins au lieu du graphe complet, et la combinaison du classifieur WBHR et de $z$ normalisé.

\begin{table}[ht]
\centering
\caption{Leviers sur l'espace fusionné, exactitude au dernier point ($\%$), moyenne de neuf runs. Entre parenthèses,
  l'écart apparié à la méthode quand il est significatif.}\label{tab-leviers}
\small
\begin{tabular}{lccc}
\toprule
configuration & HAC & esthétique & sentiment \\
\midrule
méthode & 88,64 & 65,07 & 47,35 \\
maximisation de l'information, poids 0,3 & \textbf{89,08} & 65,14 & 48,17 ($+0{,}82$, 9/9) \\
maximisation de l'information, poids 1 & 88,66 & 65,16 & \textbf{50,41} ($+3{,}07$, 7/9) \\
$\SupCon$ sur $z$ normalisé & 88,37 & 64,31 & 46,56 \\
graphe des 15 plus proches voisins & 87,48 & \textbf{65,70} & 48,39 \\
classifieur WBHR et $z$ normalisé & 88,71 & 64,36 & 47,84 \\
\midrule
plancher logistique sur $z$ (méthode) & 86,64 & 60,06 & 48,53 \\
\bottomrule
\end{tabular}
\end{table}

Aucune configuration ne change HAC ni l'esthétique au-delà du bruit. Sur le sentiment, seule la maximisation de
l'information aide ($+3{,}07$ avec le poids $1$, borne basse de l'IC à $+0{,}01$, sans correction pour comparaisons
multiples) ; elle ne rattrape ni M2CAN réglé ($-5{,}39$) ni le plancher sources seules. L'écart du graphe des plus
proches voisins à M2CAN sur l'esthétique ($+1{,}71$, IC $> 0$) n'est pas rapporté comme résultat : la configuration a
été choisie comme la meilleure des six sur le même test, et son écart à la méthode n'est pas significatif.

\subsection{Où est le goulot}

Les diagnostics localisent le facteur limitant par famille. Sur HAC, une sonde linéaire dans $z$ donne la même
exactitude que la partition et la partition est quasi parfaite : $z$ est structuré en groupes qui coïncident avec les
classes, et la géométrie les lit. Sur l'esthétique, sonde et partition coïncident aussi, à un niveau plus bas : la
limite est la séparabilité des classes dans $z$. Sur le sentiment, deux cas se présentent. Sur TumEmo, la sonde atteint
$66{,}7$ alors que la partition donne $47{,}2$ : $z$ sépare les classes mais n'a pas de structure en groupes, et la
partition échoue. Sur Yelp, la cible est isolée des sources dans $z$ (MMD$^2$ de $0{,}249$, quatre à trente fois celle
des autres cibles) et la sonde tombe à $50{,}9$ : c'est l'espace $z$ lui-même qui ne transfère pas. Les deux correctifs
qui aident sur le sentiment, la maximisation de l'information et le classifieur de WBHR, sont précisément ceux qui ne
dépendent pas de la partition.

\section{Discussion et limites}\label{sec-discussion}

\paragraph{Ce que la méthode apporte.}
Un étiquetage de la cible lu dans la géométrie de l'espace fusionné, sans classifieur, atteint le niveau de M2CAN réglé
sur deux familles sur trois, et la boucle qui renvoie cet étiquetage dans l'entraînement est le mécanisme décisif sur
HAC. La méthode est native pour $N$ sources, par le barycentre, et pour $M \ge 3$ modalités, par la fusion $M$-aire et
l'InfoNCE sur les paires de modalités, ce que M2CAN ne prévoit pas.

\paragraph{Condition d'application.}
La méthode suppose que l'espace fusionné soit structuré en groupes qui coïncident avec les classes. Le sentiment viole
cette condition et la méthode y passe sous un plancher sans adaptation, alors qu'un classifieur entraîné sur les sources
garde un pouvoir de transfert. Imposer des groupes de même taille rétablirait des partitions plausibles, mais
injecterait la distribution des classes de la cible, que l'on n'est pas censé connaître ; nous avons écarté cette
option.

\paragraph{Transductivité.}
L'inférence partitionne l'ensemble de test ensemble ; la méthode ne prédit pas un exemple isolé. Les méthodes comparées
prédisent chaque point indépendamment. Toute comparaison doit porter cette réserve.

\paragraph{Caractéristiques gelées et bout en bout.}
Tous les résultats portent sur des encodeurs gelés. Un pilote de bout en bout, encodeurs entraînés, sur AVA avec une
graine, a donné une partition effondrée (groupes de $722$ et $125$ points pour deux classes équilibrées), avec ou sans
maximisation de l'information ; les critères d'extension écrits avant le run n'étaient pas remplis et le bout en bout a
été arrêté. L'effondrement vient de $z$, non d'un déséquilibre des têtes.

\paragraph{Puissance statistique.}
Trois graines par cible et neuf paires par comparaison donnent des intervalles larges : sur HAC, un écart de
$\pm 4$ points à M2CAN n'est pas détectable. Les comparaisons multiples de la \cref{tab-leviers} ne sont pas corrigées.
Le dernier point d'évaluation évite le biais du meilleur point ; un critère d'arrêt sans étiquette fondé sur le ratio
des distances intra et inter-groupes a été essayé et perd deux points en moyenne.

\section{Conclusion}\label{sec-conclusion}

Nous avons proposé une méthode d'adaptation de domaine multi-sources et multimodale qui remplace le classifieur sur
l'espace fusionné par un étiquetage géométrique de la cible, barycentres de Wasserstein des classes sources, partition
spectrale, transport entre structures et appariement, exécuté pendant l'entraînement comme source de pseudo-étiquettes.
Dans un protocole apparié à caractéristiques gelées, elle atteint le niveau de M2CAN réglé sur la reconnaissance
d'actions et l'évaluation esthétique, et lui est inférieure sur l'analyse de sentiment. L'analyse montre que
l'appariement est optimal dans la quasi-totalité des runs, que le représentant de classe et la taille de son support
n'ont pas d'effet sur les étiquettes, et que le facteur limitant est l'existence d'une structure en groupes dans l'espace
fusionné. Les travaux futurs portent sur cette structure : combiner la maximisation de l'information et le classifieur
de WBHR, qui agissent sur des mécanismes distincts et aident tous deux sur le sentiment ; une fusion par attention,
pondérant les modalités selon l'exemple ; et une partition de la cible fondée sur des distances de Wasserstein entre
structures locales, telle que l'annonce WBHR.

\bibliographystyle{plainnat}
\bibliography{references}

\end{document}
